\chapter{熵与最大熵原理}


\section{熵(Entropy)}


假设有一个离散性随机变量$X$, 对于取值空间$R$, 其概率分布$p(x)=P(X=x)$, 有熵的定义:
\[
H(X) = - \sum_{x \in R} p(x) \log_2 p(x)
\]
约定$0 \log 0 = 0$. 在公式中, 对数的底不同对应的单位不同:
\begin{table}[H]
\centering
\begingroup
\setlength{\tabcolsep}{2pt}
\small
\begin{tabular}{|p{0.43\textwidth}|p{0.43\textwidth}|}
\hline
底数 & 单位 \\ \hline
2 & 比特(bit) \\ \hline
e & 奈特(nat) \\ \hline
10 & 哈特莱(hartley) \\ \hline
\end{tabular}
\endgroup
\end{table}
单个事件的自信息为$-\log p(x)$; 熵是自信息在分布上的期望, 可以描述一个随机变量的不确定性.
一个随机变量的熵越大, 它的不确定性越大, 那么,正确估计其值的可能性就越小.
越不确定的随机变量越需要大的信息量用以确定其值.

对于离散随机变量$X$, 熵有如下几个基本性质:

\begin{itemize}
\item 非负性: $H(X) \ge 0$.
\item 确定性变量的熵为$0$: 若$X$总是取某个固定值, 则$H(X)=0$.
\item 若$X$只有$n$个可能取值, 则$0 \le H(X) \le \log_2 n$.
\item 当$X$服从均匀分布时, 熵取得最大值$\log_2 n$.
\end{itemize}

例如, 对于一个二值随机变量$X \in \{0,1\}$, 若$P(X=1)=p$, 则
\[
H(X) = - p \log_2 p - (1-p) \log_2 (1-p)
\]
当$p=1/2$时, 熵最大, 说明二值变量在正反两种结果等可能时最不确定.

\section{最大熵模型}


最大熵原理是概率模型学习的一个准则. 最大熵原理认为, 学习概率模型时, 在所有可能的概率模型(分布)中, 熵最大的模型是最好的模型.

这一直觉可以表述为: 在已知条件之外, 不对未知信息做额外假设.
也就是说, 如果只知道一些约束条件, 那么应当选择满足这些约束且最均匀、最不偏倚的分布.

在自然语言处理的分类任务中, 常常希望学习条件概率分布$p(y|x)$. 设$f_i (x,y)$是特征函数, 经验分布为$\tilde{p}(x,y)$, 则约束通常写成:
\[
E_p [f_i] = E_{\tilde{p}}[f_i], \quad i=1,2,\dots,n
\]
其中
\[
E_{\tilde{p}}[f_i] = \sum_{x,y} \tilde{p}(x,y) f_i( x,y)
\]
\[
E_p [f_i] = \sum_{x,y} \tilde{p}(x) p(y|x) f_i (x,y)
\]

在这些约束下, 最大熵模型要求条件熵
\[
H(p) = - \sum_{x,y} \tilde{p}(x) p(y|x) \ln p(y|x)
\]
达到最大.

利用拉格朗日乘子法可以得到该模型的解具有指数形式:
\[
p(y|x) = 1/(Z(x)) \exp(\sum_i \lambda_i f_i(x,y))
\]
其中$\lambda_i$是参数, $Z(x)$是规范化因子:
\[
Z(x) = \sum_y \exp(\sum_i \lambda_i f_i(x,y))
\]

因此, 最大熵模型本质上是一个对数线性模型(log-linear model), 也可以看作指数族分布在判别建模中的应用.

最大熵模型的优点是:

\begin{itemize}
\item 能自然地容纳大量重叠特征.
\item 在约束之外保持最大的不确定性, 偏置较小.
\item 形式统一, 与逻辑回归、条件随机场等模型关系密切.
\end{itemize}

它的主要代价在于参数估计需要数值优化, 当特征很多时训练成本较高.

\section{联合熵和条件熵}


对于随机变量$X$和$Y$, 联合熵定义为:
\[
H(X,Y) = - \sum_{x,y} p(x,y) \log p(x,y)
\]
联合熵刻画的是随机变量对$(X,Y)$整体的不确定性.

条件熵定义为在已知$X$的条件下$Y$的不确定性:
\[
H(Y|X) = \sum_x p(x)H(Y|X=x)
\]
也可写成
\[
H(Y|X) = - \sum_{x,y} p(x,y) \log p(y|x)
\]

同理, 有
\[
H(X|Y) = - \sum_{x,y} p(x,y) \log p(x|y)
\]

联合熵、边缘熵与条件熵之间满足链式法则:
\[
H(X,Y) = H(X) + H(Y|X)
\]
\[
H(X,Y) = H(Y) + H(X|Y)
\]

由条件熵的定义可以看出, 条件信息会减少不确定性, 因而
\[
H(Y|X) \le H(Y)
\]
当且仅当$X$与$Y$独立时取等号; 当$Y$可由$X$唯一确定时, $H(Y|X)=0$.

\section{互信息}


互信息(mutual information)用来衡量两个随机变量共享了多少信息. 定义为:
\[
I(X;Y) = \sum_{x,y} p(x,y) \log (p(x,y)/(p(x) p(y)))
\]

它可以等价地写成以下几种形式:
\[
I(X;Y) = H(X) - H(X|Y)
\]
\[
I(X;Y) = H(Y) - H(Y|X)
\]
\[
I(X;Y) = H(X) + H(Y) - H(X,Y)
\]

因此, 互信息表示由于知道另一个变量而减少的不确定性.

互信息具有如下性质:

\begin{itemize}
\item 对称性: $I(X;Y)=I(Y;X)$.
\item 非负性: $I(X;Y) \ge 0$.
\item 当$X$与$Y$相互独立时, $I(X;Y)=0$.
\item 两个变量关联越强, 互信息通常越大.
\end{itemize}

在自然语言处理中, 互信息常用于:

\begin{itemize}
\item 衡量词与类别之间的相关程度.
\item 进行特征选择.
\item 衡量上下文与目标词之间的信息依赖关系.
\end{itemize}
